\documentclass[11pt,a4paper]{article}

% --- PACCHETTI FONDAMENTALI ---
\usepackage[utf8]{inputenc}
\usepackage[T1]{fontenc}
\usepackage[italian]{babel}
\usepackage{geometry}
\geometry{a4paper, margin=2.5cm}
\usepackage{amsmath, amssymb, amsthm}
\usepackage{graphicx}
\usepackage{booktabs}
\usepackage{hyperref}
\hypersetup{
    colorlinks=true,
    linkcolor=black,
    urlcolor=blue
}
\usepackage{xcolor}
\usepackage{listings}
\usepackage{fancyhdr}
\usepackage{caption}
\usepackage{subcaption}
\usepackage{tcolorbox} % Per definire gli ambienti promptbox e respbox

% --- COLORI E CONFIGURAZIONE CODE HIGHLIGHTING ---
\definecolor{codegreen}{rgb}{0,0.6,0}
\definecolor{codegray}{rgb}{0.5,0.5,0.5}
\definecolor{codepurple}{rgb}{0.58,0,0.82}
\definecolor{backcolour}{rgb}{0.95,0.95,0.92}
\definecolor{promptbg}{rgb}{0.94,0.97,1.0}
\definecolor{promptborder}{rgb}{0.2,0.4,0.8}
\definecolor{respbg}{rgb}{0.96,0.96,0.96}
\definecolor{respborder}{rgb}{0.5,0.5,0.5}

\newcommand{\vir}[1]{``#1''}

% Ambientazione Box personalizzate
\newtcolorbox{promptbox}[1]{
    colback=promptbg,
    colframe=promptborder,
    fonttitle=\bfseries,
    title={PROMPT #1},
    arc=2mm,
    boxrule=0.8pt
}

\newtcolorbox{respbox}[1]{
    colback=respbg,
    colframe=respborder,
    fonttitle=\bfseries,
    title={RISPOSTA #1},
    arc=2mm,
    boxrule=0.8pt
}

\lstdefinestyle{mystyle}{
    backgroundcolor=\color{backcolour},
    commentstyle=\color{codegreen},
    keywordstyle=\color{magenta},
    numberstyle=\tiny\color{codegray},
    stringstyle=\color{codepurple},
    basicstyle=\ttfamily\small,
    breakatwhitespace=false,
    breaklines=true,
    captionpos=b,
    keepspaces=true,
    numbers=left,
    numbersep=5pt,
    showspaces=false,
    showstringspaces=false,
    showtabs=false,
    tabsize=2
}
\lstset{style=mystyle}

% --- Intestazione HEADER E FOOTER ---
\pagestyle{fancy}
\fancyhf{}
\fancyhead[L]{Assignment 2}
\fancyhead[R]{Mattia Rey e Luca Ughetto}
\fancyfoot[C]{\thepage}

% --- DATI DEL DOCUMENTO ---
\title{
    \normalsize Università degli Studi di Torino \\
    \vspace{0.4cm}
    \Large \textbf{Corso di Human Computer Interaction} \\
    \vspace{0.3cm}
    \huge \textbf{Assignment 2}\\
    \Large \textbf{Gruppo DJ}
}
\author{
    \textbf{Autori:} Mattia Rey e Luca Ughetto \\
    \textbf{Matricole:}  1143424 e 1145888  \\
}
\date{2026}

\begin{document}

    \maketitle
    \thispagestyle{empty}

    \vspace{0.5cm}
    \begin{abstract}
        \noindent
        Questo documento presenta il lavoro svolto nell'ambito dell'Assignment 2 del corso di Human Computer Interaction. Esso racchiude l'analisi dettagliata e la definizione delle Personas, dei rispettivi User Journey e l'intera sequenza di prompt e risposte utilizzate per lo sviluppo iterativo del prototipo web di una piattaforma di crowdfunding.
Durante la stesura delle interazioni con il sistema di intelligenza artificiale abbiamo documentato i principali prompt che hanno posto le basi della relazione e del sito web, sottolineiamo però che per la gran parte delle modifiche abbiamo dovuto effettuare micro aggiustamenti e allineamenti al formato del sito web o della relazione. 
Questi ultimi non sono stati elencati perché di minor rilevanza ma vorremmo rendere chiaro che non ci siamo fermati alle risposte documentate facendo semplicemente un copia-incolla, abbiamo svolto diverse micro modifiche a seguito di risposte difficilmente perfette analizzando e domandandoci se, quello che ci era stato fornito, era corretto o no.

        Inoltre il sito prototipato usando Personas e User journey descritti in questa relazione è consultabile \href{https://mattiarey.github.io/HCI2026/sito/src/index.html}{qui}.
    \end{abstract}

\newpage

    \tableofcontents

\newpage

% =======================================
    \section{Definizione delle Personas}
    \label{sec:personas}
% =======================================

\subsection{Persona 1: Diego Parini (Progettista / Creator)}
\label{subsec:persona1}

\begin{table}[h!]
\centering
\small
\begin{tabular}{|p{3.5cm}|p{11.5cm}|}
\hline
\textbf{Nome} & Diego Parini \\ \hline
\textbf{Età} & 29 anni \\ \hline
\textbf{Ruolo} & Coordinatore di Progetto presso un'associazione no-profit \\ \hline
\textbf{Aggettivo} & Organizzatore \\ \hline
\textbf{Citazione} & \textit{``Voglio trasformare le idee sociali della mia comunità in progetti concreti, ma la piattaforma non deve farmi perdere tempo con errori o navigazioni confuse.''} \\ \hline
\textbf{Piattaforme d'accesso} & Desktop/Laptop (60\%), Smartphone (40\%) \\ \hline
\end{tabular}
\caption{Profilo sintetico della Persona 1}
\label{tab:persona1}
\end{table}

\subsubsection{Giornata Tipo}
\begin{itemize}
    \item \textbf{07:30 -- 08:30:} Controllo rapido di e-mail e notifiche dallo smartphone durante il tragitto sui mezzi pubblici.
    \item \textbf{09:00 -- 12:00:} Coordinamento operativo delle attività dell'associazione e gestione dei volontari.
    \item \textbf{13:00 -- 14:00:} Pausa pranzo in compagnia dei colleghi, con consultazione rapida di social media e notizie.
    \item \textbf{15:00 -- 18:00:} Gestione dei canali di comunicazione, caricamento materiali e risposte ai sostenitori.
    \item \textbf{19:00 -- 22:30:} Cena e tempo libero con la famiglia, con occasionali controlli di e-mail e notifiche.
\end{itemize}

\subsubsection{Bisogni Principali}

Raggiungere l'obiettivo economico stabilito tramite un canale affidabile e tracciabile.

Un processo chiaro e guidato per la creazione, gestione e pubblicazione della campagna.

Accesso immediato a informazioni su formazione, partner, commissioni e funzionamento del servizio.

Esperienza fluida per gestire la campagna da PC e monitorare i progressi via smartphone.

\subsubsection{Obiettivi di Miglioramento Personale}

Ridurre i tempi morti dovuti a burocrazia o problemi di usabilità dell'interfaccia.

Massimizzare la conversione delle visite in donazioni grazie a un layout chiaro ed efficace.

Sviluppare competenze pratiche sul crowdfunding e sul fundraising tramite sezioni formative integrate.

\subsubsection{Dubbi, difficoltà e possibili ostacoli}

Difficoltà nel seguire procedure complesse o poco chiare durante l'orario lavorativo.

Timore che l'interfaccia di pagamento non sia abbastanza intuitiva o sicura per i donatori meno esperti.

Paura di commettere errori irreversibili durante la fase di bozza o pubblicazione del progetto.


\subsection{Persona 2: Mirko Falletti (Sostenitore / Donatore)}
\label{subsec:persona2}

\begin{table}[h!]
\centering
\small
\begin{tabular}{|p{3.5cm}|p{11.5cm}|}
\hline
\textbf{Nome} & Mirko Falletti \\ \hline
\textbf{Età} & 65 anni \\ \hline
\textbf{Ruolo} & Ex militare ora in pensione \\ \hline
\textbf{Aggettivo} & Inflessibile \\ \hline
\textbf{Citazione} & \textit{``Vivi per il tuo paese e vivrai per sempre.''} \\ \hline
\textbf{Piattaforme d'accesso} & Desktop/Laptop (80\%), Smartphone (20\%) \\ \hline
\end{tabular}
\caption{Profilo sintetico della Persona 2}
\label{tab:persona2}
\end{table}

\subsubsection{Giornata Tipo}
\begin{itemize}
    \item \textbf{07:30 -- 08:30:} Lettura notizie e controllo e-mail dal PC fisso durante la colazione.
    \item \textbf{09:00 -- 13:00:} Tempo dedicato alla famiglia, supporto ai figli e attività con i nipoti.
    \item \textbf{14:00 -- 16:00:} Riposo pomeridiano e lettura.
    \item \textbf{18:00 -- 20:00:} Attività fisica all'aperto e cura del benessere personale.
    \item \textbf{21:00 -- 22:30:} Visione di documentari o film in compagnia della moglie e consultazione serale del PC.
\end{itemize}

\subsubsection{Bisogni Principali}

Donare e consultare i progetti senza distrazioni, pubblicità aggressive o passaggi superflui.

Testi ad alto contrasto, font ben visibili e una navigazione lineare e priva di termini troppo tecnici.

Avere la certezza assoluta di dove vadano a finire i fondi destinati alle cause supportate.

\subsubsection{Obiettivi di Miglioramento Personale}

Dedicare tempo di qualità alla famiglia recuperando gli anni trascorsi in servizio.

Contribuire attivamente a cause sociali per restituire valore alla comunità dopo l'esperienza militare.

Familiarizzare con strumenti digitali moderni senza sentirsi sopraffatto dalla tecnologia.

\subsubsection{Dubbi, difficoltà e possibili ostacoli}

Elementi grafici troppo piccoli o interfacce troppo affollate che rendono faticosa la lettura.

Esitazione nell'inserire dati bancari su piattaforme che non mostrano chiaramente standard di sicurezza.

Difficoltà nel comprendere termini inglesi o gergo tecnico non spiegato chiaramente.


\subsection{Persona 3: Alessandra Mazzoli (Attivista / Promotrice)}
\label{subsec:persona3}

\begin{table}[h!]
\centering
\small
\begin{tabular}{|p{3.5cm}|p{11.5cm}|}
\hline
\textbf{Nome} & Alessandra Mazzoli \\ \hline
\textbf{Età} & 35 anni \\ \hline
\textbf{Ruolo} & Attivista ambientale ed ecologista \\ \hline
\textbf{Aggettivo} & Empatica \\ \hline
\textbf{Citazione} & \textit{``Liberiamo i nostri mari dalla plastica, un gesto alla volta.''} \\ \hline
\textbf{Piattaforme d'accesso} & Smartphone (70\%), Desktop/Laptop (30\%) \\ \hline
\end{tabular}
\caption{Profilo sintetico della Persona 3}
\label{tab:persona3}
\end{table}

\subsubsection{Giornata Tipo}
\begin{itemize}
    \item \textbf{08:30 -- 09:30:} Colazione rapida, risveglio muscolare e controllo dei feed social legati all'ambiente.
    \item \textbf{10:00 -- 12:00:} Ricerca di campagne ecologiche in linea con i suoi valori e condivisione di contenuti sui canali social.
    \item \textbf{13:00 -- 14:00:} Pausa pranzo condivisa con altri volontari e sostenitori locali.
    \item \textbf{15:00 -- 18:00:} Sensibilizzazione sul campo nelle piazze, distribuzione di materiale informativo e condivisione di link tramite QR Code.
    \item \textbf{19:00 -- 22:30:} Cena, risposta alle interazioni online e ricerca di nuovi progetti green da sostenere o promuovere.
\end{itemize}

\subsubsection{Bisogni Principali}

Strumenti integrati per condividere istantaneamente le campagne su social media e tramite QR Code dal vivo.

Piattaforma che valorizzi cause ecologiche e progetti ad alto impatto sociale ed etico.

Accesso facilitato a informazioni su finanziamenti o iniziative dedicate allo sviluppo di prodotti eco-compatibili.

Un canale di comunicazione inclusivo che abbatta i pregiudizi e coinvolga persone di diverse estrazioni culturali.

\subsubsection{Obiettivi di Miglioramento Personale}

Ampliare la portata delle sue battaglie ambientali coinvolgendo anche i più scettici.

Promuovere l'adozione di piccoli gesti quotidiani e prodotti sostenibili a tutela dell'ecosistema.

Mantenere costanza e resilienza di fronte al disinteresse o alla lentezza dei cambiamenti sociali.

\subsubsection{Dubbi, difficoltà e possibili ostacoli}

Paura di sostenere o promuovere progetti che non siano autenticamente ecologici o trasparenti.

Difficoltà nel far accedere rapidamente le persone incontrate per strada se la pagina mobile della campagna è lenta o macchinosa.

Rischio di sentirsi sopraffatta o scoraggiata dalla mole di emergenze ambientali non affrontate adeguatamente.


\subsection{Persona 4: Luca Ponti (Sostenitore di Nicchia)}
\label{subsec:persona4}

\begin{table}[h!]
\centering
\small
\begin{tabular}{|p{3.5cm}|p{11.5cm}|}
\hline
\textbf{Nome} & Luca Ponti \\ \hline
\textbf{Età} & 22 anni \\ \hline
\textbf{Ruolo} & Studente universitario di Biotecnologie Vegetali e fotografo amatoriale \\ \hline
\textbf{Aggettivo} & Meticoloso \\ \hline
\textbf{Citazione} & \textit{``Non cerco la solita campagna di massa: voglio scovare micro-progetti dimenticati sulla biodiversità alpina prima che sia troppo tardi.''} \\ \hline
\textbf{Piattaforme d'accesso} & Smartphone (85\%), Laptop (15\%) \\ \hline
\end{tabular}
\caption{Profilo sintetico della Persona 4}
\label{tab:persona4}
\end{table}

\subsubsection{Giornata Tipo}
\begin{itemize}
    \item \textbf{06:00 -- 07:15:} Sveglia all'alba per un'uscita fotografica in collina o in serra prima delle lezioni, monitorando sul telefono notifiche da forum di settore.
    \item \textbf{08:30 -- 13:00:} Ore in laboratorio universitario e lezioni; nei minuti di pausa tra un corso e l'altro esplora piattaforme web alla ricerca di progetti scientifici dal basso.
    \item \textbf{13:30 -- 15:00:} Pranzo disordinato in aula studio mentre catalogando scatti botanici e scambia idee su gruppi Telegram dedicati al micro-fundraising etico.
    \item \textbf{16:00 -- 19:30:} Studio sul campo o sessioni di editing fotografico; test di nuove app e piattaforme esclusivamente da dispositivi mobili.
    \item \textbf{21:00 -- 23:30:} Ascolto di podcast indipendenti, ricerca di micro-cause da supportare con piccole quote mensili e messaggistica con altri appassionati.
\end{itemize}

\subsubsection{Bisogni Principali}

Possibilità di isolare campagne su tematiche di nicchia (es. tutela di specie vegetali rare, bio-archivi digitali) senza perdersi nei grandi trend.

Ricevere brevi report o diari di bordo direttamente dai creatori del progetto per vedere come viene impiegato anche un piccolo contributo.

Esperienza d'uso istantanea e senza attriti da smartphone, adatta a chi naviga in mobilità o nei momenti di pausa rapida.

Spazi di discussione o contatto diretto con i creatori per porre domande tecniche prima di effettuare la donazione.

\subsubsection{Obiettivi di Miglioramento Personale}

Creare una rete di contatti nell'ambito della ricerca scientifica indipendente ed etobotanica.

Imparare a valutare la sostenibilità reale e il rigore scientifico di un progetto prima di finanziarlo.

Documentare e dare visibilità mediatica a cause di nicchia altrimenti trascurate dai grandi media.

\subsubsection{Dubbi, difficoltà e possibili ostacoli}

Frustrazione verso piattaforme che mettono in vetrina solo cause ad alto impatto emotivo ma prive di sostanza tecnica.

Impossibilità di effettuare grandi donazioni in una sola volta, con preferenza per micro-sostegni ricorrenti senza commissioni nascoste.

Timore che le raccolte fondi non rendano conto dei risultati concreti sul lungo periodo.


\subsection{Persona 5: Isabel Marrano (Sostenitrice di Nicchia)}
\label{subsec:persona5}

\begin{table}[h!]
\centering
\small
\begin{tabular}{|p{3.5cm}|p{11.5cm}|}
\hline
\textbf{Nome} & Isabel Marrano \\ \hline
\textbf{Età} & 48 anni \\ \hline
\textbf{Ruolo} & Restauratrice di beni librari e consulente per collezioni private \\ \hline
\textbf{Aggettivo} & Esigente \\ \hline
\textbf{Citazione} & \textit{``Mi appassiona salvare le tecniche artigianali in via d'estinzione: finanzio solo chi dimostra vera perizia e rispetta le radici culturali di un territorio.''} \\ \hline
\textbf{Piattaforme d'accesso} & Desktop/Laptop (75\%), Tablet/Smartphone (25\%) \\ \hline
\end{tabular}
\caption{Profilo sintetico della Persona 5}
\label{tab:persona5}
\end{table}

\subsubsection{Giornata Tipo}
\begin{itemize}
    \item \textbf{07:00 -- 08:30:} Colazione lenta senza dispositivi digitali, seguita da una passeggiata nel quartiere e sosta nella bottega del caffè locale.
    \item \textbf{09:00 -- 13:30:} Lavoro concentrato nel laboratorio di restauro: manipolazione di materiali antichi, legature a mano e perizie tecniche senza interruzioni tecnologiche.
    \item \textbf{14:30 -- 17:00:} Appuntamenti con clienti, sopralluoghi in archivi o visite a mostre indipendenti di artigianato artistico.
    \item \textbf{17:30 -- 19:30:} Sessione al computer nello studio: risposte a e-mail di lavoro e navigazione su portali culturali e di crowdfunding d'arte.
    \item \textbf{21:00 -- 23:00:} Lettura di saggi, ascolto di musica vinile e discussione su forum specialistici su come preservare mestieri storici in pericolo.
\end{itemize}

\subsubsection{Bisogni Principali}

Schede progetto ricche di documentazione storica, foto ad alta risoluzione e spiegazioni sui metodi di lavoro adottati.

Una sezione dedicata alla conservazione del patrimonio immateriale, al recupero di vecchie manifatture e al design sostenibile di nicchia.

Preferenza per sistemi di ricompensa che offrano manufatti autentici, edizioni limitate o inviti ad workshop dedicati.

Un'interfaccia ampia, elegante e priva di notifiche invasive o layout caotici.

\subsubsection{Obiettivi di Miglioramento Personale}

Contribuire attivamente alla trasmissione dei mestieri d'arte tradizionali alle nuove generazioni.

Collezionare oggetti, pubblicazioni ed edizioni speciali ad alto valore etico e artigianale.

Utilizzare le proprie risorse finanziarie per sostenere un'economia della qualità contrapposta alla produzione di massa.

\subsubsection{Dubbi, difficoltà e possibili ostacoli}

Riluttanza a sostenere progetti che mostrano una scarsa pianificazione tecnica o un approccio puramente commerciale.

Abbandono rapido del sito se il processo di pagamento richiede troppi passaggi o se l'estetica risulta sciatta.

Timore che le ricompense promesse dai creatori vengano ritardate indefinitamente o realizzate con materiali di scarso valore.
% =======================================
    \section{User Journey delle Personas}
    \label{sec:userjourney}
% =======================================

    \subsection{Diego Parini (L'ideazione e il lancio di una campagna civica)}
    \label{subsec:uj1}

    \begin{enumerate}

        \item \textbf{Esplorazione della piattaforma e onboarding}
        \begin{itemize}
            \item \textbf{Obiettivo:} Trovare uno spazio digitale affidabile e intuitivo per lanciare la raccolta fondi del nuovo progetto di orti urbani comunitari, senza perdere ore in configurazioni ostiche.
            \item \textbf{Azioni svolte:} Navigazione dal cellulare sulla piattaforma mentre è in piedi sul tram (ore 08:00); lettura rapida della sezione ``Come funziona'' e delle FAQ sulle commissioni; apertura di una bozza di progetto inserendo titolo e prime descrizioni.
            \item \textbf{Touchpoint:} Smartphone (browser mobile della piattaforma), banner promozionali sui social media di un'altra associazione, sezione FAQ e termini di servizio del sito.
            \item \textbf{Informazioni necessarie:} Costi di gestione e trattenute, requisiti legali per aprire una campagna come no-profit, verifica della semplicità del processo di onboarding.
            \item \textbf{Pensieri ed emozioni:} \textit{``Ho pochissimo tempo prima di arrivare in ufficio, vediamo se questo sito è davvero così semplice come dicono o se devo mettermi col pc fisso stasera... Spero non chiedano diecimila documenti subito.''} (Curiosità mista a impazienza).
            \item \textbf{Difficoltà o dubbi:} Interfaccia mobile leggermente compressa nella visualizzazione delle tabelle dei costi; timore di fare passaggi falsi non potendo testare la bozza in anteprima immediata.
        \end{itemize}

        \item \textbf{Creazione e strutturazione della campagna}
        \begin{itemize}
            \item \textbf{Obiettivo:} Strutturare la campagna caricando contenuti multimediali efficaci, storie coinvolgenti e obiettivi di budget chiari per i futuri donatori.
            \item \textbf{Azioni svolte:} Passaggio al PC in ufficio (ore 15:30) per caricare foto ad alta risoluzione e video; stesura della storia del progetto e impostazione dei livelli di donazione; salvataggio della bozza e condivisione del link privato con il board dell'associazione.
            \item \textbf{Touchpoint:} Dashboard di creazione campagna da Desktop, servizio di cloud storage locale per recuperare le immagini, chat interna con il team di lavoro.
            \item \textbf{Informazioni necessarie:} Suggerimenti di scrittura persuasiva (storytelling) integrati, formati e limiti di peso per i file multimediali, anteprima fedele della pagina su PC e smartphone.
            \item \textbf{Pensieri ed emozioni:} \textit{``Il testo rende bene, ma ho paura che il video sia troppo pesante e non si carichi. Voglio che la pagina ispiri fiducia immediata a chi arriva dai nostri social.''} (Concentrazione, senso di responsabilità).
            \item \textbf{Difficoltà o dubbi:} Mancanza di un feedback visivo immediato sulla formattazione del testo descrittivo; dubbio su come strutturare i 'tier' di donazione per renderli appetibili.
        \end{itemize}

        \item \textbf{Revisione e pubblicazione del progetto}
        \begin{itemize}
            \item \textbf{Obiettivo:} Verificare l'affidabilità della piattaforma e completare la validazione formale prima di schiacciare il tasto di pubblicazione definitiva.
            \item \textbf{Azioni svolte:} Invio della richiesta di revisione al team di moderazione; lettura delle linee guida sulla sicurezza dei pagamenti da mostrare ai sostenitori; controllo dello stato della pratica tramite e-mail.
            \item \textbf{Touchpoint:} Casella di posta elettronica, sezione centro assistenza della piattaforma, pagina di riepilogo pre-pubblicazione.
            \item \textbf{Informazioni necessarie:} Tempistiche certe di approvazione della campagna (es. ``entro 24 ore''), contatti diretti con un operatore in caso di intoppi burocratici.
            \item \textbf{Pensieri ed emozioni:} \textit{``Spero che i controlli non vadano per le lunghe, abbiamo pianificato la conferenza stampa per giovedì e il link deve essere attivo per forza!''} (Ansia da scadenza, trepidazione).
            \item \textbf{Difficoltà o dubbi:} Sensazione di 'scatola nera' durante l'attesa dell'approvazione senza un tracciamento chiaro in tempo reale.
        \end{itemize}

        \item \textbf{Monitoraggio e interazione con la community}
        \begin{itemize}
            \item \textbf{Obiettivo:} Monitorare l'andamento delle prime donazioni e ringraziare la community direttamente in mobilità.
            \item \textbf{Azioni svolte:} Apertura dell'app mobile durante la pausa pranzo (ore 13:30); visualizzazione delle notifiche di donazione; invio di un messaggio di aggiornamento e ringraziamento alla bacheca pubblica.
            \item \textbf{Touchpoint:} Smartphone (notifiche push o web-app mobile), pannello di controllo ``Creator'' semplificato.
            \item \textbf{Informazioni necessarie:} Statistiche in tempo reale (totale raccolto, numero donatori), strumenti rapidi per pubblicare aggiornamenti sui social direttamente dalla piattaforma.
            \item \textbf{Pensieri ed emozioni:} \textit{``Fantastico, le prime notifiche stanno arrivando dai vicini di quartiere! Devo pubblicare subito un grazie per mantenere alta l'energia della campagna.''} (Entusiasmo, gratificazione).
            \item \textbf{Difficoltà o dubbi:} Difficoltà nel rispondere singolarmente ai donatori se la rubrica messaggi mobile è poco responsiva.
        \end{itemize}

    \end{enumerate}

    \vspace{0.3cm}
    \noindent\textbf{Opportunità di Progetto Emergenti:} Dalla sintesi del percorso di Diego emergono le seguenti necessità progettuali per la piattaforma: un onboarding rapido e guidato (wizard) con calcolatore dei costi integrato; l'adozione di \textit{Smart Templates} e suggerimenti contestuali per la stesura della campagna; un tracker di moderazione visivo in tempo reale per monitorare lo stato di approvazione; e una mini-dashboard Creator ottimizzata per dispositivi mobili che permetta interazioni rapide (ringraziamenti e condivisioni) in mobilità.


    \subsection{Mirko Falletti (Il sostenitore cauto e concreto)}
    \label{subsec:uj2}

    \begin{enumerate}

        \item \textbf{Ricerca del progetto e verifica dell'affidabilità}
        \begin{itemize}
            \item \textbf{Obiettivo:} Trovare una causa locale o nazionale seria, a favore di famiglie di militari o comunità in difficoltà, da sostenere in modo trasparente e diretto.
            \item \textbf{Azioni svolte:} Accensione del PC fisso dopo cena (ore 21:00); ricerca tramite Google o clic su un link condiviso via e-mail; navigazione nella homepage alla ricerca di una categoria chiara.
            \item \textbf{Touchpoint:} Desktop/Laptop, link ricevuto via e-mail o forum di settore, menu di navigazione della piattaforma.
            \item \textbf{Informazioni necessarie:} Garanzie ufficiali sull'affidabilità dell'ente promotore, indicazione immediata dell'impatto e della percentuale di fondi destinata alla causa, layout con caratteri nitidi.
            \item \textbf{Pensieri ed emozioni:} \textit{``Spero che non sia uno di quei siti moderni pieni di video che partono da soli o scritte minuscole... Voglio solo capire chi c'è dietro questo progetto.''} (Diffidenza iniziale, prudenza).
            \item \textbf{Difficoltà o dubbi:} Disturbo visivo di fronte a banner animati o pop-up, incertezza sulla reale legittimità dell'associazione promotrice.
        \end{itemize}

        \item \textbf{Analisi della campagna e scelta dell'importo}
        \begin{itemize}
            \item \textbf{Obiettivo:} Analizzare in dettaglio la scheda della campagna selezionata per verificare come verranno impiegati i fondi prima di decidere l'importo.
            \item \textbf{Azioni svolte:} Lettura approfondita del testo e controllo del prospetto di spesa; verifica dei documenti allegati; selezione dell'opzione di donazione singola, rifiutando abbonamenti automatici.
            \item \textbf{Touchpoint:} Scheda dettaglio della campagna (sezione ``Trasparenza''), modale/pulsante ``Dona Ora'' con scelta dell'importo.
            \item \textbf{Informazioni necessarie:} Ripartizione chiara (grafico a torta o tabella) di dove va ogni euro donato, assenza di caselle pre-spuntate per mancie non richieste.
            \item \textbf{Pensieri ed emozioni:} \textit{``Vedo che hanno specificato esattamente quanti mattoni acquisteranno. Bene. Non voglio sorprese sul conto corrente il mese prossimo.''} (Soddisfazione, bisogno di controllo).
            \item \textbf{Difficoltà o dubbi:} Paura che scatti un abbonamento mensile non voluto o che le commissioni per la piattaforma siano ambigue e nascoste.
        \end{itemize}

        \item \textbf{Pagamento online e ricezione della ricevuta}
        \begin{itemize}
            \item \textbf{Obiettivo:} Effettuare il pagamento online in modo sicuro e senza intoppi tecnici o procedure di autenticazione troppo astruse.
            \item \textbf{Azioni svolte:} Inserimento dei dati bancari; conferma della transazione dall'app della propria banca sullo smartphone; attesa del messaggio a schermo.
            \item \textbf{Touchpoint:} Form di pagamento sicuro, app bancaria personale su smartphone, schermata di ringraziamento.
            \item \textbf{Informazioni necessarie:} Indicazione esplicita della sicurezza (lucchetto SSL), ricevuta scaricabile immediatamente in formato PDF da stampare e archiviare.
            \item \textbf{Pensieri ed emozioni:} \textit{``Ecco, adesso mi chiede il codice di verifica sullo smartphone... Spero di fare in tempo prima che scada la sessione.''} (Lieve tensione da prestazione tecnologica).
            \item \textbf{Difficoltà o dubbi:} Disorientamento se la pagina reindirizza a un sito esterno non brandizzato; mancanza di una ricevuta fiscale chiara per il commercialista.
        \end{itemize}

        \item \textbf{Ricezione degli aggiornamenti post-donazione}
        \begin{itemize}
            \item \textbf{Obiettivo:} Verificare periodicamente che i fondi abbiano prodotto il risultato promesso, senza dover accedere ai social media.
            \item \textbf{Azioni svolte:} Apertura dell'e-mail di riepilogo mensile; lettura del report fotografico sui lavori; salvataggio dell'aggiornamento sul PC.
            \item \textbf{Touchpoint:} E-mail periodica di aggiornamento (Newsletter sostenitori), link al report di impatto sul sito.
            \item \textbf{Informazioni necessarie:} Resoconto sintetico dei risultati, nessun tono spinto verso nuove donazioni, ma pura rendicontazione etica.
            \item \textbf{Pensieri ed emozioni:} \textit{``Mi hanno inviato la foto della struttura ultimata direttamente via e-mail. È così che si lavora, la serietà prima di tutto.''} (Orgoglio, fiducia consolidata).
            \item \textbf{Difficoltà o dubbi:} Rischio che gli aggiornamenti finiscano nello Spam o che richiedano log-in complessi per essere letti.
        \end{itemize}

    \end{enumerate}

    \vspace{0.3cm}
    \noindent\textbf{Opportunità di Progetto Emergenti:} Il percorso di Mirko evidenzia l'importanza della chiarezza e della sicurezza visiva. Si suggerisce l'implementazione di: Badge di Verifica visibili per gli enti promotori; modalità di navigazione ad ``Alta Leggibilità''; prospetti trasparenti per la ripartizione del budget; un check-out lineare senza spunte preselezionate per abbonamenti; un generatore istantaneo di ricevute PDF ai fini fiscali; e l'invio di aggiornamenti ``No-Login'' fruibili direttamente via e-mail.


    \subsection{Alessandra Mazzoli (L'attivista dinamica e promotrice di rete)}
    \label{subsec:uj3}

    \begin{enumerate}

        \item \textbf{Esplorazione e verifica dell'impatto etico}
        \begin{itemize}
            \item \textbf{Obiettivo:} Scovare progetti ecologici ad alto impatto etico che necessitano di visibilità e verificare che non si tratti di un'operazione di greenwashing.
            \item \textbf{Azioni svolte:} Navigazione veloce da smartphone durante la colazione (ore 08:45); filtraggio campagne nella categoria ``Ambiente e Sostenibilità''; controllo delle certificazioni della filiera.
            \item \textbf{Touchpoint:} Smartphone (web app responsive), link condivisi nei gruppi Telegram di attivismo, sezione ``Certificazioni ed Etica'' del progetto.
            \item \textbf{Informazioni necessarie:} Dettagli trasparenti sull'impatto ecologico (es. kg di plastica rimossi), assenza di partnership controverse, materiale informativo condivisibile.
            \item \textbf{Pensieri ed emozioni:} \textit{``Questa causa contro la plastica nei mari ha un potenziale enorme, ma devo assicurarmi che sia un progetto vero prima di metterci la faccia.''} (Entusiasmo cauto, spirito critico).
            \item \textbf{Difficoltà o dubbi:} Mancanza di dati chiari sulla tracciabilità delle azioni; difficoltà nel distinguere cause a reale impatto sociale da mere campagne commerciali.
        \end{itemize}

        \item \textbf{Adesione e generazione del kit di promozione}
        \begin{itemize}
            \item \textbf{Obiettivo:} Aderire alla campagna con un contributo e ottenere strumenti digitali o materiali per ingaggiare la propria rete sia online che di persona.
            \item \textbf{Azioni svolte:} Donazione rapida via smartphone (es. Apple Pay/Satispay); generazione automatica del link di referral e del QR Code personalizzato; download delle grafiche per le stories social.
            \item \textbf{Touchpoint:} Modale di pagamento istantaneo su mobile, dashboard ``Ambassador'', generatore automatico di QR Code e asset grafici.
            \item \textbf{Informazioni necessarie:} Disponibilità di QR Code ad alta risoluzione pronti da mostrare dallo schermo, testi brevi ed efficaci (pitch) pronti per le caption dei social.
            \item \textbf{Pensieri ed emozioni:} \textit{``Fantastico! Se ho il QR Code pronto sul telefono posso farlo scansione subito ai ragazzi durante la sensibilizzazione in piazza oggi pomeriggio!''} (Energia, senso di efficacia).
            \item \textbf{Difficoltà o dubbi:} Link di condivisione troppo lunghi o assenza di formati grafici adatti a Instagram/TikTok; donazione macchinosa che frena le conversioni a freddo.
        \end{itemize}

        \item \textbf{Sensibilizzazione e condivisione dal vivo}
        \begin{itemize}
            \item \textbf{Obiettivo:} Coinvolgere le persone incontrate sul campo (piazze, eventi, banchetti) facilitando l'accesso immediato alla campagna.
            \item \textbf{Azioni svolte:} Attività nelle piazze (ore 16:30); mostra il QR Code dal proprio smartphone ai passanti; li guida a compilare i pochi campi necessari per una firma o micro-donazione.
            \item \textbf{Touchpoint:} Smartphone della persona incontrata per strada (landing page d'atterraggio da QR Code), interazione humana diretta.
            \item \textbf{Informazioni necessarie:} Caricamento ultrarapido della landing page anche con connessione mobile lenta (4G/3G), messaggio sintetico nei primi 3 secondi di schermata.
            \item \textbf{Pensieri ed emozioni:} \textit{``Spero che la pagina si carichi subito, la gente in strada ha fretta e se il sito ci mette 5 secondi a caricare perdo l'attenzione del passante.''} (Adrenalina, determinazione).
            \item \textbf{Difficoltà o dubbi:} Pagine pesanti che falliscono il caricamento su reti instabili; moduli di registrazione troppo lunghi che scoraggiano l'utente occasionale.
        \end{itemize}

        \item \textbf{Monitoraggio dell'impatto di rete}
        \begin{itemize}
            \item \textbf{Obiettivo:} Verificare i risultati del proprio lavoro di sensibilizzazione e mantenere ingaggiata la community creata.
            \item \textbf{Azioni svolte:} Consultazione serale (ore 21:30) delle metriche del proprio link/QR Code; pubblicazione di un ringraziamento sui social per il traguardo raggiunto; partecipazione alle discussioni in bacheca.
            \item \textbf{Touchpoint:} Dashboard del promotore (statistiche di rete), canali social personali, bacheca aggiornamenti della piattaforma.
            \item \textbf{Informazioni necessarie:} Numero di persone ingaggiate tramite il proprio QR Code, totale dei fondi raccolti grazie alla condivisione diretta.
            \item \textbf{Pensieri ed emozioni:} \textit{``Vedere che grazie ai miei scatti e ai QR in piazza siamo riusciti a raccogliere altri 500 euro mi dà una carica pazzesca per la prossima azione!''} (Orgoglio, senso di appartenenza).
            \item \textbf{Difficoltà o dubbi:} Impossibilità di vedere il proprio impatto di rete indiretto (es. le donazioni generate dai link re-condivisi da un amico).
        \end{itemize}

    \end{enumerate}

    \vspace{0.3cm}
    \noindent\textbf{Opportunità di Progetto Emergenti:} L'esperienza di Alessandra spinge per un'integrazione fluida tra mondo fisico e digitale: è necessario introdurre un Bollino ``Eco-Trust'' con relativo \textit{Impact Index} certificato; strutturare un ``Social \& Street Ambassador Kit'' che generi istantaneamente QR Code e materiale per i social; sviluppare landing page ultra-leggere pensate per conversioni immediate in strada con connettività ridotta; e implementare un sistema di gamification etica che mostri graficamente l'impatto di rete generato (es. ``Grazie a te, 12 persone hanno donato'').


    \subsection{Luca Ponti (Il sostenitore meticoloso di micro-cause scientifiche)}
    \label{subsec:uj4}

    \begin{enumerate}

        \item \textbf{Ricerca di micro-cause scientifiche}
        \begin{itemize}
            \item \textbf{Obiettivo:} Scovare micro-progetti di conservazione botanica o archivi digitali sulla biodiversità alpina, ignorati dalle grandi raccolte fondi di massa.
            \item \textbf{Azioni svolte:} Esplorazione rapida da smartphone durante la pausa universitaria (ore 13:15); utilizzo di filtri per parole chiave iper-specifiche; salvataggio tra i preferiti dei progetti rigorosi.
            \item \textbf{Touchpoint:} Smartphone (app o web app), motore di ricerca interno con filtri di nicchia, link segnalati su gruppi Telegram di appassionati.
            \item \textbf{Informazioni necessarie:} Sistema di tag dettagliato (non solo la macro-categoria ``Ambiente''), schede sintetiche con premesse metodologiche e scientifiche, identificazione chiara del team di ricerca.
            \item \textbf{Pensieri ed emozioni:} \textit{``Spero di non dovere scorrere cinquanta campagne generiche sui gattini prima di trovare un progetto di tutela locale... Serve un filtro per la ricerca seria.''} (Scetticismo verso il mainstream).
            \item \textbf{Difficoltà o dubbi:} Tendenza della piattaforma a mostrare in prima pagina solo le campagne ad alto impatto emotivo nascondendo quelle a basso budget ma di alto valore scientifico.
        \end{itemize}

        \item \textbf{Valutazione tecnica e dialogo con i ricercatori}
        \begin{itemize}
            \item \textbf{Obiettivo:} Analizzare il rigore e la sostenibilità tecnica del progetto scelto prima di vincolare le proprie limitate risorse da studente.
            \item \textbf{Azioni svolte:} Esame approfondito della scheda e dei PDF allegati; invio di una domanda diretta tramite la chat al team per chiarire la metodologia; lettura delle credenziali accademiche.
            \item \textbf{Touchpoint:} Scheda dettagliata del progetto da smartphone, modulo Q\&A diretto con i creatori, allegati tecnici in PDF.
            \item \textbf{Informazioni necessarie:} Metriche di successo verificabili, trasparenza sui costi delle attrezzature e delle uscite sul campo, risposte competenti dai creatori.
            \item \textbf{Pensieri ed emozioni:} \textit{``Il progetto è ottimo, ma voglio capire come intendono conservare i semi raccolti. Invio un messaggio al team: se rispondono bene, attivo una donazione mensile.''} (Rigore analitico).
            \item \textbf{Difficoltà o dubbi:} Assenza di uno spazio pubblico per porre domande tecniche senza inviare mail formali; descrizioni limitate a storytelling emotivi privi di dati.
        \end{itemize}

        \item \textbf{Attivazione del sostegno ricorrente}
        \begin{itemize}
            \item \textbf{Obiettivo:} Impostare un sostegno finanziario ricorrente di piccolo taglio (micro-donazione) in linea con le proprie limitate disponibilità economiche.
            \item \textbf{Azioni svolte:} Selezione della micro-donazione mensile (es. 5€); pagamento rapido con portafoglio digitale da smartphone (Apple Pay/Google Pay); attivazione notifiche per il ``diario di bordo''.
            \item \textbf{Touchpoint:} Selettore di importo con opzione ricorrente, gateway di pagamento mobile, schermata di conferma.
            \item \textbf{Informazioni necessarie:} Garanzia di totale assenza di commissioni a sorpresa; massima facilità di disdetta o pausa dell'abbonamento con un solo tap.
            \item \textbf{Pensieri ed emozioni:} \textit{``Non posso dare 100€ in una volta, ma 5€ al mese per un anno sono fattibili. L'importante è che possa disattivare l'addebito se ho un mese stretto.''} (Pragmatismo economico).
            \item \textbf{Difficoltà o dubbi:} Paura di rimanere vincolato a contratti difficili da cancellare; importi minimi prefissati troppo elevati per un micro-sostegno.
        \end{itemize}

        \item \textbf{Consultazione dei dati e condivisione accademica}
        \begin{itemize}
            \item \textbf{Obiettivo:} Seguire l'evoluzione dei rilevamenti botanici e condividere gli avanzamenti scientifici nei propri cerchi accademici.
            \item \textbf{Azioni svolte:} Ricezione di notifiche push con i resoconti periodici; consultazione delle nuove foto e dei dati grezzi pubblicati; condivisione su forum di settore o gruppi universitari.
            \item \textbf{Touchpoint:} Notifiche push o e-mail dedicate, sezione ``Diario di Bordo'' del progetto, pulsante rapido di condivisione su gruppi Telegram/Discord.
            \item \textbf{Informazioni necessarie:} Aggiornamenti autentici e non patinati dal campo (diari, foto, dati grezzi), possibilità di scaricare campioni di dati.
            \item \textbf{Pensieri ed emozioni:} \textit{``Hanno appena pubblicato le coordinate fotografiche della specie tutelata... Vedere i dati veri che arrivano dalla montagna vale ogni centesimo versato.''} (Entusiasmo scientifico).
            \item \textbf{Difficoltà o dubbi:} Aggiornamenti generici inviati raramente e privi di reale contenuto di approfondimento scientifico.
        \end{itemize}

    \end{enumerate}

    \vspace{0.3cm}
    \noindent\textbf{Opportunità di Progetto Emergenti:} Il profilo di Luca suggerisce lo sviluppo di un'infrastruttura adatta alla ricerca di nicchia: filtri di ricerca verticali e iper-specifici uniti a una sezione ``Hidden Gems''; l'introduzione di una tab tecnica integrata con un modulo Peer-to-Peer Q\&A; l'implementazione di sistemi di \textit{Micro-Supporter Subscriptions} flessibili e disdicibili con un tap; e un feed di aggiornamento in stile ``Field Log'' che supporti la pubblicazione di open data e rilevamenti sul campo anziché solo materiale promozionale.


    \subsection{Isabel Marrano (La sostenitrice esigente di beni culturali e artigianato)}
    \label{subsec:uj5}

    \begin{enumerate}

        \item \textbf{Individuazione di progetti di restauro e artigianato}
        \begin{itemize}
            \item \textbf{Obiettivo:} Individuare progetti di recupero di antichi mestieri d'arte o di restauro di beni librari/artistici in cerca di finanziamenti etici.
            \item \textbf{Azioni svolte:} Navigazione da computer fisso nel tardo pomeriggio (ore 18:00); filtraggio nelle categorie ``Patrimonio Culturale'' e ``Design Sostenibile''; selezione di schede con un'estetica curata e foto ad alta risoluzione.
            \item \textbf{Touchpoint:} Desktop/Laptop con schermo ampio, portali di notizie d'arte o newsletter culturali, galleria dei progetti in evidenza sulla piattaforma.
            \item \textbf{Informazioni necessarie:} Presentazione estetica impeccabile senza elementi grafici caotici, dettagli sulle radici storiche della bottega, immagini di grande formato che mostrino la lavorazione manuale.
            \item \textbf{Pensieri ed emozioni:} \textit{``Voglio capire se dietro questa campagna per il recupero della legatoria storica c'è vera maestria o solo una buona operazione di marketing... L'immagine è bella, approfondiamo.''} (Curiosità selettiva).
            \item \textbf{Difficoltà o dubbi:} Fastidio di fronte a layout sciatti, testi troppo piccoli o moduli confusamente organizzati; sospetto di scarsa professionalità se mancano fonti storiche.
        \end{itemize}

        \item \textbf{Valutazione delle ricompense e dei materiali}
        \begin{itemize}
            \item \textbf{Obiettivo:} Valutare il progetto e selezionare un livello di donazione che offra un'opera o una ricompensa tangibile ad alto valore etico.
            \item \textbf{Azioni svolte:} Lettura approfondita della storia del laboratorio e delle biografie dei maestri; esame delle ricompense offerte (es. volumi rilegati a mano); confronto tra i diversi livelli di supporto.
            \item \textbf{Touchpoint:} Scheda estesa del progetto (tab sulla storia e sul processo creativo), modulo dei livelli di ricompensa (Tier di supporto) su desktop.
            \item \textbf{Informazioni necessarie:} Descrizione dettagliata dei materiali utilizzati (carte pregiate, inchiostri naturali), stima realistica dei tempi di produzione, trasparenza sui costi di produzione rispetto alla quota destinata alla bottega.
            \item \textbf{Pensieri ed emozioni:} \textit{``Questa edizione in 50 esemplari numerati è uno splendido esempio di artigianato. So di finanziare la trasmissione di un mestiere in via di estinzione.''} (Appagamento culturale).
            \item \textbf{Difficoltà o dubbi:} Incertezza sulla qualità effettiva dei manufatti, timore di ritardi indefiniti nella consegna o mancanza di dettagli sull'origine dei materiali impiegati.
        \end{itemize}

        \item \textbf{Checkout ed emissione dell'ordine}
        \begin{itemize}
            \item \textbf{Obiettivo:} Finalizzare il contributo importante tramite un processo di pagamento sicuro, elegante e privo di intoppi.
            \item \textbf{Azioni svolte:} Selezione della quota di donazione con ricompensa (es. 150€); inserimento dell'indirizzo di spedizione e dei dati della carta su PC; revisione del riepilogo e invio dell'ordine.
            \item \textbf{Touchpoint:} Carrello e pagina di checkout per desktop, form di inserimento dati di fatturazione, e-mail di conferma ordine.
            \item \textbf{Informazioni necessarie:} Garanzia sulla sicurezza delle transazioni, tutela dei dati personali, ricevuta dettagliata con l'indicazione specifica della ricompensa prenotata.
            \item \textbf{Pensieri ed emozioni:} \textit{``Il processo è stato fluido e composto. Apprezzo che il sito non mi chieda di condividere l'acquisto su tre social network diversi appena ho finito di pagare.''} (Rassicurazione e apprezzamento per la sobrietà).
            \item \textbf{Difficoltà o dubbi:} Richieste invasive di condivisione social obbligatoria o pop-up rumorosi alla fine della transazione che rovinerebbero l’esperienza.
        \end{itemize}

        \item \textbf{Tracciamento della produzione e ricezione}
        \begin{itemize}
            \item \textbf{Obiettivo:} Tracciare lo stato di lavorazione del manufatto prenotato e fruire del report finale del progetto di restauro.
            \item \textbf{Azioni svolte:} Consultazione della newsletter mensile del laboratorio con foto delle fasi di lavorazione; riconferma dell'indirizzo all'avvicinarsi della spedizione; ricezione e verifica del materiale prodotto.
            \item \textbf{Touchpoint:} E-mail di aggiornamento sulla produzione (``Dietro le quinte''), notifica di tracciamento del corriere, oggetto fisico con packaging etico.
            \item \textbf{Informazioni necessarie:} Aggiornamenti puntuali in caso di fisiologici ritardi nella produzione artigianale, codice di tracciamento chiaro per la consegna.
            \item \textbf{Pensieri ed emozioni:} \textit{``Il volume è arrivato ed è di una fattura straordinaria. Questa piattaforma seleziona davvero eccellenze.''} (Fidelizzazione, ammirazione, senso di appartenenza).
            \item \textbf{Difficoltà o dubbi:} Silenzio prolungato da parte del creatore durante i lunghi mesi necessari alla lavorazione manuale.
        \end{itemize}

    \end{enumerate}

    \vspace{0.3cm}
    \noindent\textbf{Opportunità di Progetto Emergenti:} Per soddisfare i requisiti estetici e qualitativi di utenti come Isabel, la piattaforma dovrebbe fornire: un layout di tipo ``Editorial \& Heritage'' che valorizzi immagini ad altissima risoluzione e testi storici; una sezione per schede ricompensa trasparenti arricchite con tracciabilità delle materie prime e certificati digitali di autenticità; un flusso di checkout elegante e sobrio (senza spinte alla condivisione forzata); e un \textit{Craftsmanship Progress Tracker} visivo che mantenga aggiornato il sostenitore sulle fasi della lavorazione artigianale (es. reperimento materiali $\rightarrow$ lavorazione $\rightarrow$ controllo qualità).
% =======================================
    \section{Appendice: Prompt e Interazione AI}
    \label{sec:appendice}
% =======================================

    Questo capitolo raccoglie la sequenza di prompt forniti all'Intelligenza Artificiale e le relative risposte ottenute durante le fasi di progettazione e sviluppo del sito di crowdfunding (GINGER Lab), coprendo l'analisi delle Personas, degli User Journey e lo sviluppo iterativo del prototipo web.

    \subsection{Prompt 1: Creazione della prima Persona (Diego Parini)}
    \label{prompt:1}
    \begin{promptbox}{1}
        Devo fare un progetto, scritto in latex, per un corso di human computer interaction.\\
        Dobbiamo creare personas e user journey per creare un sito basato sugli utenti che lo andranno ad utilizzare. Questo sito si occupa di crowdfunding.\\
        Iniziamo creando una prima Personas. Quindi inventati una persona fittizia con: nome, ruolo, aggettivo rappresentativo, una sua citazione, cosa fa in una sua giornata tipo, quali sono i suoi bisogni più importanti, come vorrebbe migliorare la propria vita, cosa ostacola il raggiungimento dei suoi obiettivi e quali sono le piattaforme che usa per accedere al sito.\\
        \vspace{0.2cm}
        \textit{Il pdf che ti incollo è l'analisi di un sito simile che non funziona}
    \end{promptbox}

    \begin{respbox}{1}
        Profilo generato per la Persona 1 (Diego Parini), inclusivo di ruolo, citazione, giornata tipo, bisogni principali e ostacoli nell'uso della piattaforma di crowdfunding.
    \end{respbox}

    \subsection{Prompt 2: Criteri dettagliati per le Personas (Mirko Falletti e Alessandra Mazzoli)}
    \label{prompt:2}
    \begin{promptbox}{2}
        Devo fare un progetto, scritto in latex, per un corso di human computer interaction.\\
        Dobbiamo creare personas e user journey per creare un sito basato sugli utenti che lo andranno ad utilizzare. Questo sito si occupa di crowdfunding.\\
        Ogni persona deve includere:\\
        - nome e breve descrizione del profilo;\\
        - contesto personale o professionale rilevante per l’uso del sito;\\
        - obiettivi e motivazioni;\\
        - bisogni informativi;\\
        - competenze digitali e familiarità con il crowdfunding;\\
        - dubbi, difficoltà e possibili ostacoli;\\
        - elementi che possono favorire o ridurre la fiducia nell’iniziativa;\\
        - uno scenario d’uso concreto.\\
        Potete considerare, ad esempio, persone interessate a sostenere la ricerca, membri della comunità universitaria, persone alla prima esperienza di crowdfunding oppure ricercatrici e ricercatori interessati a proporre una campagna. Scegliete e motivate i profili più pertinenti al vostro progetto.
    \end{promptbox}

    \begin{respbox}{2}
        Dettagli completi per la Persona 2 (Mirko Falletti) e per la Persona 3 (Alessandra Mazzoli), focalizzati sul contesto d'uso, bisogni informativi e fattori di fiducia.
    \end{respbox}

    \subsection{Prompt 3: Creazione di altre due Personas (Luca Ponti e Isabel Marrano)}
    \label{prompt:3}
    \begin{promptbox}{3}
        Creiamo altre due personas diverse dalle precedenti, ti do un po' di informazioni e poi crea tu: sii creativo, non devono sembrare dei robot creati da un IA, usa informazioni non classiche! Età: 22 e 48, nomi: Luca Ponti e Isabel Marrano, vogliono sostenere dei progetti che gli stanno a cuore su tematiche non ancora trattate e abbastanza di nicchia. Per quanto riguarda la giornata tipo, anche lì, sii creativo e varia molto
    \end{promptbox}

    \begin{respbox}{3}
        Elaborazione delle Personas 4 e 5 (Luca Ponti e Isabel Marrano) con focus su sostenitori di nicchia, approcci alla ricerca personalizzati ed esigenze estetiche/metodologiche specifiche.
    \end{respbox}

    \subsection{Prompt 4: Unione di tutte e 5 le Personas in un unico file LaTeX}
    \label{prompt:4}
    \begin{promptbox}{4}
        Unisci tutte e 5 le Personas in un unico file LaTeX completo e ben formattato.
    \end{promptbox}

    \begin{respbox}{4}
        Generazione del blocco unico LaTeX formattato con tabelle e liste strutturate per le 5 Personas.
    \end{respbox}

    \subsection{Prompt 5: Creazione del primo User Journey (Diego Parini)}
    \label{prompt:5}
    \begin{promptbox}{5}
        Ora basandoti sulla personas 1, crea uno user journey inserendo:\\
        - obiettivo della persona;\\
        - azioni svolte;\\
        - touchpoint, cioè punti di contatto con il sito o con altri canali;\\
        - informazioni necessarie;\\
        - pensieri ed emozioni ipotizzati;\\
        - difficoltà o dubbi;\\
        - opportunità di miglioramento da tradurre nel progetto. Mi raccomando, agisci come prima, sii creativo e non aver paura di uscire dagli schemi con l'inventiva
    \end{promptbox}

    \begin{respbox}{5}
        Generazione del primo User Journey (Diego Parini) in forma tabellare con fasi del percorso, touchpoint e opportunità d'interfaccia.
    \end{respbox}

    \subsection{Prompt 6: Creazione sequenziale degli User Journey (da Persona 2 a 5)}
    \label{prompt:6}
    \begin{promptbox}{6}
        Ora partendo dalla personas 2, una alla volta fino alle 5, crea uno user journey inserendo:\\
        - obiettivo della persona;\\
        - azioni svolte;\\
        - touchpoint, cioè punti di contatto con il sito o con altri canali;\\
        - informazioni necessarie;\\
        - pensieri ed emozioni ipotizzati;\\
        - difficoltà o dubbi;\\
        - opportunità di miglioramento da tradurre nel progetto. Mi raccomando, agisci come prima, sii creativo e non aver paura di uscire dagli schemi con l'inventiva. Al termine di questo messaggio, quando ti scriverò continua, vai alla personas successiva
    \end{promptbox}

    \begin{respbox}{6}
        Sviluppo sequenziale degli User Journey per le Personas 2, 3, 4 e 5 con mappatura dettagliata dell'esperienza utente.
    \end{respbox}

    \subsection{Prompt 7: Progettazione del primo prototipo di sito web (8 pagine)}
    \label{prompt:7}
    \begin{promptbox}{7}
        Dobbiamo realizzare un primo prototipo di sito web (HTML e CSS) basandoci sulle personas e gli user journey. Si tratta di un sito di crowdfunding e dobbiamo ispirarci ad un sito precedentemente analizzato, nel PDF allegato è presente proprio codesta analisi. Progetta il sito in ampiezza, definendo le principali sezioni, le voci di menu e la struttura di navigazione. Sviluppa poi in profondità almeno un percorso completo e significativo, collegato a una delle personas.\\
        È richiesta la versione desktop. Devi creare almeno 8 pagine del sito con contenuti o funzioni differenti tra loro
    \end{promptbox}

    \begin{respbox}{7}
        Fornite le prime 8 pagine web con i file CSS incorporati al loro interno.
    \end{respbox}

    \subsection{Prompt 8: Separazione del CSS in file dedicati}
    \label{prompt:8}
    \begin{promptbox}{8}
        Forniscimi le diverse pagine web HTML ma dividi il CSS in uno o più file dedicati!
    \end{promptbox}

    \begin{respbox}{8}
        Ristrutturazione del codice con la separazione pulita tra file HTML e stylesheet CSS dedicati.
    \end{respbox}

    \subsection{Prompt 9: Correzione della ricerca e aggiunta di progetti dinamici}
    \label{prompt:9}
    \begin{promptbox}{9}
        C'è un problema nella ricerca di cause, progetti, partner. La ricerca effettivamente non funziona! Inoltre nella pagina esplora progetti, vediamo solo il singolo progetto presente, vorrei che rendessi più dinamica questa parte aggiungendo progetti inventati. Anche aggiungendo (solo se necessario) altre parti di codice o memorizzazione
    \end{promptbox}

    \begin{respbox}{9}
        Implementazione del file JavaScript per la memorizzazione e la visualizzazione dinamica di un set esteso di progetti.
    \end{respbox}

    \subsection{Prompt 10: Risoluzione dei bug su cerca.html}
    \label{prompt:10}
    \begin{promptbox}{10}
        problemi rilevati: cerca.html ha ancora problemi, se inserisco qualunque carattere, anche riconducibile ad un effettivo progetto spariscono tutti e ancora, se rimuovo il mio input, non ricompaioni! Inoltre il filtraggio non funziona
    \end{promptbox}

    \begin{respbox}{10}
        Risoluzione completa dei bug su \texttt{cerca.html}: riprogettazione della gestione degli eventi in JavaScript, ripristino istantaneo degli elementi al cancellare dell'input e corretto funzionamento del filtraggio per categorie e tag.
    \end{respbox}

    \subsection{Prompt 11: Correzione del nome del progetto nella pagina checkout.html}
    \label{prompt:11}
    \begin{promptbox}{11}
        Problemi rilevati: Nella pagina checkout.html la scritta "Stai sostenendo: " indica sempre lo stesso progetto.
    \end{promptbox}

    \begin{respbox}{11}
        Modifica del codice JavaScript per la pagina checkout.html in modo che la scritta "Stai sostenendo: " rifletta correttamente il progetto selezionato dall'utente, utilizzando variabili dinamiche e aggiornamento del DOM.
    \end{respbox}

    \subsection{Prompt 12: Inserimento dell'indice all'inizio del file}
    \label{prompt:12}
    \begin{promptbox}{12}
        Analizza il contenuto di questo file LaTeX e aggiungi nella parte inziale l'indice mancante. Regole da seguire: Mantieni l'indentazione utilizzata in tutto il file già presente nel risultato prodotto. Non modificare nulla al di fuori della table of contents. Restituiscimi la nuova parte generata da inserire nel file 
    \end{promptbox}

    \begin{respbox}{12}
        Modifica del codice LaTeX della relazione per inserire l'indice correlato al contenuto di essa
    \end{respbox}

    \subsection{Prompt 13: Connessione dei singoli prompt}
    \label{prompt:13}
    \begin{promptbox}{13}
	Inserisci i tag mancanti nel resto del file per rendere funzionanti i link dell'indice e restituiscimi il file LaTeX finale. Regole: Non modificare l'indentazione fornita. Modifica solo lo stretto necessario che ti ho indicato.
    \end{promptbox}

    \begin{respbox}{13}
        Integrazione della funzionalità di connessione e referenziazione puntuale dei singoli prompt all'interno dell'indice generale e della documentazione.
    \end{respbox}


\newpage

\end{document}